\subsection{Inverse images of manifold structures, submanifolds}

5.8.1. Let X be a topological space, Y be a manifold, and f be a map from X to Y. Consider the following conditions:

(QR) (resp. (R)) For every \(a \in X\) , there exists an open neighborhood \(U\) of a in \(X\) and a chart \((V, \varphi, E)\) of the manifold \(Y\) at \(f(a)\) such that \(f(U) \subset V\) and that \(\varphi \circ f\) induces a homeomorphism of \(U\) onto the intersection of \(\varphi(V)\) with a closed vector subspace (resp. a closed vector subspace admitting a topological complement) \(F\) of \(E\) .

If condition (QR) is satisfied, there exists on X one and only one manifold structure for which the triples (U, \(\varphi \circ f\) , F) (with the notations of (QR)) are charts of X. It is called the inverse image structure of the manifold structure of Y by \(f\) . Its topology is that of X.

For there to exist on X a manifold structure compatible with the given topology and for which \(f\) is an immersion, it is necessary and sufficient that condition (R) be satisfied. This structure is then unique: it is the inverse image under f of the manifold structure of Y.

5.8.2. The condition (R) above is satisfied in particular when for every \(a \in X\) there exists an open neighborhood U of \(a\) such that \(f|U\) is a homeomorphism from \(U\) onto an open subset of \(Y\). In this case, the manifold \(X\) obtained is étale over \(Y\) (5.7.6).

5.8.3. Suppose now that X is a topological subspace of Y, with f being the canonical injection. If f satisfies condition (R) (resp. (QR)) of no. 5.8.1, then X, equipped with the inverse image structure of the manifold structure of Y under f, is said to be a submanifold (resp. quasi-submanifold) of Y. A submanifold is a quasi-submanifold.

Every quasi-submanifold is locally closed; every open set is a submanifold and its structure is that defined in No. 5.2.3.

When Y is locally of finite dimension, a topological subspace X of Y is a submanifold of Y if and only if, for every \(a \in X\) there exists an open neighborhood U of a in Y, a coordinate system \((\xi^{1}, \ldots, \xi^{m})\) of Y in U, and an integer \(k \leqslant m\) such that \(U \cap X\) is the set of points of U where the functions \(\xi^{i}\) for \(1 \leqslant i \leqslant k\) vanish simultaneously.

Let \(f:X\to Y\) be a morphism of varieties. Suppose that \(f\) is an immersion, and let \(f\) induce a homeomorphism from X onto \(f(X)\) . Then \(f(X)\) is a submanifold of Y, and \(f\) induces an isomorphism from X onto \(f(X)\) ; one then says that \(f\) is an embedding of X into Y.

5.8.4. Let Y and W be two varieties, let X be a subvariety of Y, and let f be a map from X to W. For f to be a morphism, it is necessary and sufficient that, for every \(a \in X\) , there exist a neighbourhood U of a in Y and a morphism \(g : U \to W\) such that g coincides with f on \(U \cap X\) .

5.8.5. Let X be a quasi-subvariety of a variety Y, and let g be a map from a variety Z into X. Suppose that K is of characteristic 0, or that X is a subvariety of Y. For g to be a morphism from Z to X, it is necessary and sufficient that it be a morphism from Z to Y: in other words, the variety structure of X is induced by that of Y (Ens., ch.~IV, § 2, no. 4).

5.8.6. If X is a subvariety (resp. quasi-subvariety) of Y and if Y is a subvariety of a variety Z, then X is a subvariety (resp. quasi-subvariety) of Z.

5.8.7. Let \(X_{i}\) (for i = 1, 2) be a variety and let \(Y_{i}\) be a subvariety (resp. quasi-subvariety) of \(X_{i}\) (for i = 1, 2). Then \(Y_{1} \times Y_{2}\) is a subvariety (resp. quasi-subvariety) of \(X_{1} \times X_{2}\) .

5.8.8. Let X be a quasi-subvariety of a variety Y; let x be a point of X and denote by f the canonical injection of X into Y. The map \(T_{x}(f):T_{x}(X)\to T_{x}(Y)\) is injective and allows one to identify the space \(T_{x}(X)\) with a closed vector subspace of \(T_{x}(Y)\) . The quotient topological vector space \(T_{x}(Y)/T_{x}(X)\) is a Banach space, which is called the transversal space to X at x (in Y). Its dimension (finite or +∞) is called the codimension of X in Y at the point x.

If moreover X is a subvariety of Y, then the space \(T_{x}(X)\) admits a topological supplement in \(T_{x}(Y)\) .

5.8.9. Let \(f\) be a map from a variety X to a variety Y, and let \(\Gamma\) be its graph. For \(f\) to be a morphism, it is necessary and sufficient that the following two conditions be satisfied:

\begin{enumerate}
\def\labelenumi{(\roman{enumi})}
\item
  \(\Gamma\) is a subvariety of \(X \times Y\) .
\item
  For every \((x, y) \in \Gamma\) , one has
\end{enumerate}

\[
\mathrm{T} _ {(x, y)} (\mathrm{X} \times \mathrm{Y}) = \mathrm{T} _ {(x, y)} (\Gamma) \oplus \mathrm{T} _ {y} (\mathrm{Y}).
\]

If this is so, the map \(\mathsf{pr}_1\) induces an isomorphism from \(\Gamma\) onto \(X\) , and \(\mathrm{T}_{(x,y)}(\Gamma)\) is identified with the graph of \(\mathrm{T}_x(f)\) .

In particular, the diagonal of \(X \times X\) is a subvariety of \(X \times X\) .

5.8.10. Let Y be a variety, and let \((f_{i})_{i\in I}\) be a finite family of morphic functions on Y. Let X be the set of \(x\in Y\) such that \(f_{i}(x)=0\) for every i. Assume the following hypothesis:

\begin{enumerate}
\def\labelenumi{(\Alph{enumi})}
\setcounter{enumi}{9}
\tightlist
\item
  For every \(x \in X\) , the differentials \(d_x f_i\) are linearly independent in \(T_x'(Y)\) .
\end{enumerate}

Then X is a closed subvariety of Y, the tangent space \(T_{x}(X)\) is the subspace of \(T_{x}(Y)\) formed by the \(\alpha\) such that \(\alpha.f_{i}=0\) for every i in I. Moreover, the codimension of X in Y is equal to Card (I) at each of its points.

5.8.11. (Simple points of an ideal). Let \(\mathfrak{a}\) be an ideal of the polynomial algebra \(\mathrm{K}[\mathrm{X}_1,\ldots ,\mathrm{X}_n]\) . A point \(x = (x_{1},\dots,x_{n})\) of \(\mathbf{K}^n\) is called a zero of \(\mathfrak{a}\) if \(f(x) = 0\) for all \(f\in \mathfrak{a}\) . If \(x\in \mathbf{K}^n\) , denote by \(\mathbf{S}_x\) the subalgebra of \(\mathrm{K}(\mathrm{X}_1,\dots,\mathrm{X}_n)\) formed by the fractions \(f / g\) , with \(f,g\in \mathrm{K}[\mathrm{X}_1,\dots,\mathrm{X}_n]\) , and \(g(x)\neq 0\) ; denote by \(\mathfrak{a}_x\) the ideal of \(\mathbf{S}_x\) generated by \(\mathfrak{a}\) in \(\mathbf{S}_x\) . A point \(x\) is called a simple zero of \(\mathfrak{a}\) if it is a zero of \(\mathfrak{a}\) and if the following condition is satisfied:

\begin{enumerate}
\def\labelenumi{(\Alph{enumi})}
\setcounter{enumi}{18}
\tightlist
\item
  There exists a finite sequence \((f_1, \ldots, f_m)\) of elements of \(\mathfrak{a}\) which generate the ideal \(\mathfrak{a}_x\) and whose differentials at \(x\) are linearly independent. (This condition is equivalent to saying that the local ring \(\mathbf{S}_x / \mathfrak{a}_x\) is regular (Alg. Comm., to appear).)
\end{enumerate}

Let \(Z\) (resp. \(Z_s\) ) be the set of zeros (resp. simple zeros) of \(\mathfrak{a}\). The set \(Z\) is closed in \(K^n\) , \(Z_s\) is open in \(Z\) , and \(Z_s\) is a subvariety of \(K^{n}\) . If \(x \in Z_{s}\) , the ideal \(\mathrm{K}[\mathrm{X}_{1}, \ldots, \mathrm{X}_{n}] \cap \mathfrak{a}_{x}\) consists of the polynomials f vanishing on a neighborhood of x in \(Z_{s}\) .

Let \(\overline{\mathfrak{a}}\) be the ideal of polynomials vanishing on Z. If K is algebraically closed, the set of simple zeros of \(\mathfrak{a}\) is dense in Z.

5.8.12. Let X be a variety and L be the set of pairs \((x, Z)\) where x is a point of X and Z a subvariety of X containing x. Given two pairs \(\pi = (x, Z)\) and \(\pi' = (x', Z')\) , we denote by \(R\{\pi, \pi'\}\) the relation:

"We have \(x = x'\) and there exists a neighborhood U of \(x\) such that \(\mathbf{U} \cap \mathbf{Z} = \mathbf{U} \cap \mathbf{Z'}\) . Then R is an equivalence relation on L; we denote by \(\gamma_x(Z)\) the equivalence class of the pair \((x, Z) \in L\) . On the set \(J = L/R\) , there exists one and only one variety structure satisfying the following condition:

For every subvariety Z of X, the map \(x \mapsto \gamma_{x}(Z)\) from Z into J is an isomorphism of Z onto an open subvariety of J.

We say that J is the variety of germs of subvarieties of X (cf.~Top. Gén., chap.~I, 4 \(^{e}\) ed., § 6, no. 10).

The map \(\rho: J \to X\) defined by \(\rho(\gamma_x(Z)) = x\) is an immersion; it is called the canonical immersion of \(J\) to \(X\) .

If X is a separated analytic variety of finite dimension, so is J.

